\section{Introduction}

\label{sec:introduction}

Agent memory systems face a fundamental tension between recency and relevance. A memory item stored months ago may contain the exact fact needed for a query today; an item stored seconds ago may share only surface vocabulary without containing its answer. Retrieval selectors that treat freshness as a primary ordering signal risk displacing relevant older items with irrelevant newer ones. We call this failure mode temporal dominance, and we show that it is the structural cause of a large and consistent accuracy deficit across every evaluated token budget in the system we study.

The problem is precise and algebraic. The deployed hybrid scoring function assigns fixed weights to lexical Jaccard overlap, temporal decay, access frequency, and heuristic importance. Because the temporal weight is query-independent, a freshness gap of moderate size can outweigh a positive lexical advantage whenever Jaccard scores are small -- which is the typical case when a short natural-language question is compared against long conversational memory segments. We derive the exact condition under which displacement occurs and show it is routinely satisfied in the evaluated bank.

We propose CTLS (Contrastive Tier-Based Lexical Selection) as the principled corrective. CTLS partitions candidates into a lexically relevant tier (Jaccard overlap J(q,m) strictly positive) and a lexically absent tier (J(q,m) equal to zero) before applying any temporal scoring. Within the relevant tier, the full hybrid score ranks candidates so that temporal decay and frequency serve as tiebreakers among genuinely relevant items. Packing draws from the relevant tier first, then from the absent tier if budget permits. This partition provably satisfies a Non-Dominance Property: no zero-overlap item can outrank a positive-overlap item regardless of freshness gap, frequency, or importance values. Temporal dominance is eliminated by construction, not by parameter tuning.

Our empirical evidence comes from a controlled evaluation on a stratified LongMemEval-S subset of 84 memory episodes. All selectors draw from the same frozen candidate bank with identical token budgets, metadata, memory wrapper, and whole-item first-fit packing. Only the scoring function changes. We evaluate across three budgets -- 512, 1024, and 2048 tokens -- with 3 answering calls per memory episode, stratified paired cluster bootstrap confidence intervals, and Holm correction across all prespecified comparisons.

BM25 outperforms the deployed hybrid at every budget by -36.1111 percentage points, -29.3651 percentage points, and -33.7302 percentage points (all Holm-corrected). Removing only the temporal term recovers most of this deficit, confirming it as the primary cause. CTLS is the principled generalization: where hybrid\_no\_time discards temporal information globally, CTLS preserves it as a within-tier tiebreaker. The CTLS direct experimental evaluation remains future work; this paper establishes the theoretical guarantee and the empirical motivation.

\subsection{Quality under Matched Budgets}

\label{sec:rq_quality}

RQ1 asks whether CTLS closes the gap between the deployed hybrid and BM25 under matched token budgets on the frozen bank. BM25 outperforms the hybrid by -36.1111 percentage points, -29.3651 percentage points, and -33.7302 percentage points at the three evaluated budgets (all Holm-corrected). CTLS is designed to close this gap by eliminating temporal dominance through tier partition while preserving temporal decay as a within-tier tiebreaker rather than discarding it globally.

\subsection{Mechanism Attribution via Tier Partition}

\label{sec:rq_attribution}

RQ2 asks which component of the hybrid scoring function causes the deficit and whether the tier partition addresses it by construction. The mechanism analysis derives the temporal dominance condition: the temporal weight can outweigh a positive Jaccard advantage whenever the freshness gap exceeds a threshold set by the weight ratio and the half-life parameter. Three empirical consequences follow: the deficit should be largest in single-session-user strata; removing only the temporal term should recover most of the loss; and evidence hit rates should fall with accuracy rates. All three are visible in the data. CTLS formalizes the fix while preserving within-tier temporal ordering among genuinely relevant items.

\subsection{Lifecycle Cost under Memory Reuse}

\label{sec:rq_lifecycle}

RQ3 asks how write amortization changes per-query cost when a single memory bank serves multiple distinct queries, and whether CTLS changes that cost. CTLS is a zero-cost selector substitution: no model calls, no embedding model, no parameters beyond the existing Jaccard tokenizer. Construction cost, write amortization, and per-query token counts are identical between CTLS and any selector on the same bank. At 160 distinct queries against one instance, the total token ratio of model-acknowledgement to raw write paths is 1.21.
